\section{Condition-Aware Inference}

\label{sec:condition}

\subsection{Image Quality Estimation}

Field photographs exhibit systematic quality degradation compared to curated datasets.
The CAI module estimates image quality along four axes using lightweight (sub-5ms) estimators:

\begin{equation}
    \mathbf{q}(\mathbf{x}) = (q_{\text{blur}}, q_{\text{exposure}}, q_{\text{occlusion}}, q_{\text{distance}})
\end{equation}

Each estimator is a small convolutional network (3 layers, 50K parameters) trained on the ZooBench-Field dataset with human-annotated quality labels.

\subsection{Adaptive Preprocessing}

Based on the quality estimate, CAI selects from a library of preprocessing configurations:

\begin{table}[H]
\centering
\caption{Condition-adaptive preprocessing pipeline.}
\label{tab:preprocessing}
\begin{tabularx}{\textwidth}{@{}llX@{}}
\toprule
\textbf{Condition} & \textbf{Preprocessing} & \textbf{Rationale} \\
\midrule
Low light & Histogram equalization + denoising & Recover detail from underexposed images \\
Motion blur & Multi-scale inference (3 crops) & Capture subject at multiple blur scales \\
Partial occlusion & Sliding window attention & Focus on visible body parts \\
Long distance & Super-resolution (2x) + center crop & Enhance small subject in frame \\
High quality & Standard resize only & Avoid unnecessary processing overhead \\
\bottomrule
\end{tabularx}
\end{table}

\subsection{Confidence Calibration}

Field conditions systematically affect model confidence calibration.
CAI applies condition-specific temperature scaling:

\begin{equation}
    p_{\text{cal}}(s \mid \mathbf{x}) = \text{softmax}\left(\frac{\mathbf{z}(\mathbf{x})}{T(\mathbf{q}(\mathbf{x}))}\right)
\end{equation}

\noindent where $T(\mathbf{q})$ is a temperature function learned on the ZooBench-Field validation set, mapping quality vectors to calibration temperatures.
Under degraded conditions, the temperature increases, producing more uniform (less overconfident) probability distributions.

When calibrated confidence falls below a threshold, CAI provides guidance rather than a definitive identification:

\begin{equation}
    \text{Output}(\mathbf{x}) = \begin{cases}
        \text{ID}(s^*) & \text{if } p_{\text{cal}}(s^* \mid \mathbf{x}) \geq \tau_{\text{high}} \\
        \text{TopK}(s_1, \ldots, s_k) & \text{if } \tau_{\text{low}} \leq p_{\text{cal}}(s^* \mid \mathbf{x}) < \tau_{\text{high}} \\
        \text{Guidance}(\mathbf{q}) & \text{if } p_{\text{cal}}(s^* \mid \mathbf{x}) < \tau_{\text{low}}
    \end{cases}
\end{equation}

\noindent where $\text{Guidance}(\mathbf{q})$ suggests specific actions to improve image quality (e.g., ``Move closer to the animal'' or ``Turn on flash'').

% ============================================================
% 6. ZooBench-Field Benchmark
% ============================================================
